% Chapter 5, Topic _Linear Algebra_ Jim Hefferon
%  http://joshua.smcvt.edu/linearalgebra
%  2016-Jun-17
% \documentclass{article}
% \usepackage{graphicx}
% \usepackage{amsmath}

% \newcommand{\definend}[1]{\textit{#1}}
\topic{耦合振子}
\index{Coupled Oscillators|(}
% \begin{document}

这是一个\definend{Wilberforce 摆}。\index{Wilberforce pendulum}
弹簧上悬挂着一个质量，也就是摆锤。
将它向上推一点然后松手，它就会上下振动。
\begin{center}
  \includegraphics{jc/asy/wilber.pdf}
\end{center}
而且，如果把这个装置调节得当，就会出现十分迷人的现象。
几秒钟后，除了上下运动之外，
质量开始绕着
沿弹簧中线竖直向上的轴转动。
这种偏航转动不断增大，直到运动几乎完全变成旋转，
而上下运动所剩无几。
大约五秒之后，运动又演化为两者的混合。
再过一段时间，规律性重新出现。
令人惊奇的是，此时
运动几乎完全是竖直的。
这一过程持续进行，装置交替经历
纯竖直运动的阶段与纯转动运动的阶段，
其间还夹杂着混合状态。
（在网上搜索 ``wilberforce pendulum video''，
可以看到一些精彩的演示。）

每一个纯运动状态都是振动的一个\definend{简正模}。
我们将分析该装置处于简正模时的行为。
这全都与特征值有关。

\begin{center}
  \includegraphics{jc/asy/wilber001.pdf}
\end{center}

用~$x(t)$ 表示随时间变化的竖直运动，
$\theta(t)$ 表示转动运动。
取定坐标系，使得在静止位置 $x=0$ 且~$\theta=0$，
并使得正的 $x$ 方向朝上，且正的 $\theta$
从上方看去为逆时针方向。

我们先来对一个弹簧上的质量的运动建模，
此时弹簧被限制为不发生扭转。
这比较简单，因为只有一种运动，即一个自由度。
将质量置于静止位置，
再向上推动它以压缩弹簧。
胡克定律说的是，当距离较小时，
回复力
与距离成正比，即 $F=-k\cdot x$。
常数~$k$ 是弹簧的
\definend{劲度}。

牛顿定律指出，力与相应的加速度成正比，
$F=m\cdot d^2\,x(t)/dt$。
比例常数 $m$ 是这个物体
的\definend{质量}。
将胡克定律与牛顿定律结合起来，就得到描述
该质量运动的微分方程
$m\cdot d^2\,x(t)/dt=-k\cdot x(t)$。
我们更喜欢把变量都写在一边的形式。
\begin{equation*}
  m\cdot \frac{d^2\,x(t)}{dt}+k\cdot x(t)=0
  \tag{$*$}
\end{equation*}

我们的物理直觉是，摆锤会随时间振荡。
它从高处出发，因此
图像应当如下。
\begin{center}
  \includegraphics{jc/asy/wilber003.pdf}
\end{center}
当然，这看起来像余弦函数的图像，我们也认识到，
运动微分方程~($*$) 被
$x(t)=\cos \omega t$ 满足，其中 $\omega=\sqrt{m/k}$，
这是因为 $dx/dt=\omega\cdot \sin \omega t$
且 $d^2x/dt^2=-\omega^2\cdot \cos \omega t$。
这里，$\omega$ 是\definend{角频率}。
它决定振荡的周期，因为若 $\omega=1$ 则周期为~$2\pi$，
而若 $\omega=2$ 则周期为~$\pi$，等等。

我们可以给出~($*$) 的更一般的解。
对于一般的振幅，我们在前面添上一个因子~$A$，
$x(t)=A\cdot \cos wt$。
并且还可以允许有相位移动，
这样，我们就不必在质量处于最高点时开始计时，
而是取
$x(t)=A\cos (wt+\phi)$。
这就是\definend{简谐运动}的方程。

现在回到对耦合的一对运动的讨论，
即竖直运动与转动。
这两种运动之所以相互作用，是因为被扭转的弹簧会
稍稍伸长或缩短；
例如，它可能像这里这样起作用，也可能反过来，这取决于
弹簧。
\begin{center}
  \includegraphics{jc/asy/wilber002.pdf}
\end{center}
反过来，一个从静止位置被拉伸或压缩的弹簧
也会稍稍扭转。
这两者的相互作用产生\definend{耦合振荡}。

要理解这种相互作用如何产生我们在简正模中
看到的引人注目的行为，
可以设想质量正朝着会使弹簧变长的方向
旋转。
如果与此同时竖直方向的运动是弹簧
正在变短，
那么把两者叠加起来
就可能使它们几乎相互抵消。
结果摆锤几乎不再做竖直运动，只是扭转。
装置若调节得当，这种状态可以持续数秒之久。

所谓``调节得当''，是指
纯竖直运动的周期
与纯转动运动的周期相同或接近。
这样，抵消就会持续一段时间。

运动之间的相互作用也能产生另一种简正模
行为：摆锤几乎只做竖直运动而转动不多，
这发生在弹簧的运动~$x(t)$
产生的扭转与摆锤的扭转~$\theta(t)$ 相抗衡时。
此时摆锤几乎停止转动，
其运动几乎完全是竖直的。

要得到这个两个自由度情形下的运动方程，
我们做一个与一个自由度情形相同的假设：
对于小的位移，回复力
与位移成正比。
但现在我们把这一假设同时用于
竖直运动与转动。
设转动运动中的比例常数为~$\kappa$。
类似地，对转动运动我们也使用力与加速度
成正比的牛顿定律，并取
比例常数为~$I$。

最关键的是，我们在这两种运动之间添加一个耦合项，
它与每一个运动都成正比，比例
常数为~$\epsilon/2$。

这就得到一个由两个微分方程构成的方程组，
第一个描述竖直运动，第二个
描述转动。
这些方程描述耦合系统在任意时刻~$t$ 的行为。
\begin{equation*}
  \begin{split}
  m\cdot\frac{d^2\,x(t)}{dt^2}
      +k\cdot x(t)+\frac{\epsilon}{2}\cdot \theta(t) 
      &=0  \\
  I\cdot\frac{d^2\,\theta(t)}{dt^2}
      +\kappa\cdot \theta(t)+\frac{\epsilon}{2}\cdot x(t)  
      &=0
  \end{split}
  \tag{$**$}
\end{equation*}
我们将用它们来分析
该系统处于简正模时的行为。

首先考虑不耦合的运动，即不含
$\epsilon$ 项的方程所给出的运动。
没有这些项时，它们描述的是简谐函数，
我们记~$\omega_x^2$ 为~$k/m$，$\omega_\theta^2$ 为 $\kappa/I$。
前面已经论证过，要观察到这种停驻行为，
应当把装置调节得使两个周期相同，即 $\omega_x^2=\omega_\theta^2$。
记~$\omega_0$ 为这个数。

现在考虑耦合的运动 $x(t)$ 与~$\theta(t)$。
根据同样的道理，要观察到停驻行为，我们要两者
同步，例如，当伸长达到峰值时转动也恰好处于
峰值。
也就是说，在简正模中，两个振荡具有相同的
角频率~$\omega$。
至于相位移动，正如前面也曾讨论的，有两种
情形：弹簧的运动所施加的扭转与转动振荡给出的扭转
方向相同，或者两者方向相反。
无论哪种情形，要得到简正模，峰值都必须重合。
\begin{align*}
  x(t) &= A_1\cos(\omega t+\phi)  
       &x(t) &= A_1\cos(\omega t+\phi)   \\  
  \theta(t) &= A_2\cos(\omega t+\phi)  
       &\theta(t) &= A_2\cos(\omega t+(\phi+\pi))    
\end{align*}
我们将仔细推演左边的情形，把另一种情形留作练习。

我们要找出哪些 $\omega$ 是可能的。
取二阶导数
\begin{equation*}
  \frac{d^2\,x(t)}{dt}=-A_1\omega^2\cos(\omega t+\phi)
  \qquad
  \frac{d^2\,\theta(t)}{dt}=-A_2\omega^2\cos(\omega t+\phi)
\end{equation*}
并代入运动方程组~($**$)。
\begin{align*}
  m\cdot(-A_1\omega^2\cos(\omega t+\phi))
      +k\cdot (A_1\cos(\omega t+\phi))
      +\frac{\epsilon}{2}\cdot (A_2\cos(\omega t+\phi)) 
      &=0  \\
  I\cdot(-A_2\omega^2\cos(\omega t+\phi))
      +\kappa\cdot (A_2\cos(\omega t+\phi))
      +\frac{\epsilon}{2}\cdot (A_1\cos(\omega t+\phi))  
      &=0  
\end{align*}
提出 $\cos(\omega t+\phi)$，再用~$m$ 去除。
\begin{align*}
  \big(\frac{k}{m}-\omega^2\big)\cdot A_1 
      +\frac{\epsilon}{2m}\cdot A_2 
      &=0  \\
  \big(\frac{\kappa}{I}-\omega^2\big)\cdot A_2
      +\frac{\epsilon}{2m}\cdot A_1  
      &=0  
\end{align*}
我们假定 $k/m=\omega_{x}^2$
与 $\kappa/I=\omega_{\theta}^2$ 相等，
并把该数记为 $\omega_0^2$。
作此替换，并将其改写为矩阵方程。
\begin{equation*}
  \begin{mat}
    \omega_0^2-\omega^2    &\epsilon/2m  \\
    \epsilon/2I &\omega_0^2-\omega^2 
  \end{mat}
  \colvec{
    A_1 \\
    A_2}
  =
  \begin{mat}
    0 \\
    0
  \end{mat}
\end{equation*}
显然，该方程组有平凡解~$A_1=0$，$A_2=0$，
对应质量静止不动的情形。
我们想知道，对于哪些频率~$\omega$，该方程组有非平凡
解。
\begin{equation*}
  \begin{mat}
    \omega_0^2    &\epsilon/2m  \\
    \epsilon/2I &\omega_0^2
  \end{mat}
  \colvec{
    A_1 \\
    A_2}
  =
  \omega^2
  \begin{mat}
    A_1 \\
    A_2
  \end{mat}
\end{equation*}
简正模的角频率~$\omega$ 就是该矩阵的特征值。

要求它，
取行列式并令其为零。
\begin{equation*}
  \begin{vmatrix}
    \omega_0^2-\omega^2    &\epsilon/2m  \\
    \epsilon/2I &\omega_0^2-\omega^2
  \end{vmatrix}=0 \quad\Longrightarrow\quad
  \omega^4-(2\omega_0^2)\,\omega^2 +(\omega_0^4
  -\frac{\epsilon^2}{4mI})=0
\end{equation*}
该方程是关于~$\omega^2$ 的二次方程。
应用求解二次方程的公式 $(-b\pm\sqrt{b^2-4ac})/(2a)$。
\begin{equation*}   % \rule[-.3\baselineskip]{0pt}{\baselineskip}
  \omega^2 
  =\frac{2\omega_0^2
    \pm\sqrt{\smash[b]{\strut} 4\omega_0^4
      -4(\omega_0^4-\epsilon^2/4mI)}}{2}
  =\omega_0^2\pm\frac{\epsilon}{2\sqrt{mI}}
\end{equation*}
值 $\epsilon/\sqrt{mI}=\epsilon/\sqrt{\kappa k}$ 常记作
$\omega_B$，于是
$
\omega^2 =\omega_0^2\pm\omega_B/2
$。
这就是\definend{拍频}，即两个简正模频率
之差。

尽管这一论证超出了本书的范围，但摆的运动的一般公式
是各简正模中的运动
的线性组合。
因此，摆的运动完全由上述矩阵的特征值
决定。
参见 \cite{BergMarshall}。

\begin{exercises}
\item 利用和的余弦公式，
  给出简谐运动的一个更一般的公式。
  \begin{answer}
    余弦函数的和角公式为
    $\cos(\alpha+\beta)=cos(\alpha)\cos(\beta)-\sin(\alpha)\sin(\beta)$。
    把 $A\cos(\omega t+\phi)$ 展开为 $A\cdot[cos(\omega
    t)\cos(\phi)-\sin(\omega t)\sin(\phi)]$。由于 $\cos(\phi)$ 与
    $\sin(\phi)$ 不随~$t$ 变化，所以我们得到通解
    $x(t)=B\cos(\omega t)+C\sin(\omega t)$。
  \end{answer}
\item 求与这些特征值相对应的特征向量。
  \begin{answer}
    我们对第一个根进行计算；第二个类似。
    我们有如下方程。
    \begin{equation*}
      \begin{mat}
        \omega_0^2-\omega^2  &\epsilon/2m \\
        \epsilon/2I         &\omega_0^2-\omega^2
      \end{mat}
      \colvec{A_1 \\ A_2}
      =\colvec{0 \\ 0}
    \end{equation*}
    代入第一个根~$\omega^2=\omega_0^2+\epsilon/2\sqrt{mI}$。
    两个方程中有一个是多余的，所以我们只考虑第一个。
    \begin{equation*}
      -(\epsilon/2\sqrt{mI})\cdot A_1+(\epsilon/2m)\cdot A_2=0
    \end{equation*}
    这是平面中一条过原点的直线，因此要确定它，
    我们只需求出第一个变量与第二个变量之比。
    \begin{equation*}
      \frac{A_2}{A_1}=\frac{(\epsilon/2\sqrt{mI})}{(\epsilon/2m)}
                     =\sqrt{\frac{m}{I}}
    \end{equation*}
    所以这就是与第一个特征值相对应的特征向量空间。
    \begin{equation*}
      \set{\colvec{A_1 \\ \sqrt{m/I}\cdot A_1}\suchthat A_1\in\C}
    \end{equation*}
  \end{answer}
\item 求在如下情形中 $\omega$ 的取值：
  $x(t)=A_1\cos(\omega t+\phi)$
  且 $\theta(t)=A_2\cos(\omega t+(\phi+\pi))$。
  \begin{answer}
    注意到 $\cos(\omega t+(\phi+\pi))=-\cos(\omega t+\phi)$。
    取导数
    \begin{equation*}
      \frac{d\,x(t)}{dt}=-\omega\cdot\sin(\omega t+\phi)
      \qquad
      \frac{d^2\,x(t)}{dt^2}=-\omega^2\cdot\cos(\omega t+\phi)
    \end{equation*}
    以及
    \begin{equation*}
      \frac{d\,\theta(t)}{dt}=\omega\cdot\sin(\omega t+\phi)
      \qquad
      \frac{d^2\,\theta(t)}{dt^2}=\omega^2\cdot\cos(\omega t+\phi)
    \end{equation*}
    并代入运动方程组~($**$)。
    \begin{align*}
      mA_1\cdot(-\omega^2\cos(\omega t+\phi))
          +kA_1\cdot (\cos(\omega t+\phi))
          +\frac{\epsilon}{2}A_2\cdot (-\cos(\omega t+\phi)) 
          &=0  \\
      IA_2\cdot(\omega^2\cos(\omega t+\phi))
          +\kappa A_2\cdot (-\cos(\omega t+\phi))
          +\frac{\epsilon}{2}A_1\cdot (\cos(\omega t+\phi))  
          &=0  
    \end{align*}
    提出 $\cos(\omega t+\phi)$
    \begin{align*}
     \big(m(-\omega^2)\big)\cdot A_1
         +kA_1 
         -\frac{\epsilon}{2}\cdot A_2 
         &=0  \\
     \big(I\omega^2\big)\cdot A_2
         -\kappa A_2
         +\frac{\epsilon}{2}\cdot A_1  
         &=0  
   \end{align*}
   再用~$m$ 和~$I$ 去除。
   \begin{align*}
    \big(\frac{k}{m}-\omega^2\big)\cdot A_1 
        -\frac{\epsilon}{2m}\cdot A_2 
        &=0  \\
    \frac{\epsilon}{2I}\cdot A_1  
    -\big(\frac{\kappa}{I}-\omega^2\big)\cdot A_2
        &=0  
  \end{align*}
   我们假定 $k/m=\omega_{x}^2$
   与 $\kappa/I=\omega_{\theta}^2$ 相等，
   并把该数记为 $\omega_0^2$。
   作此替换，并将其改写为矩阵方程。
   \begin{equation*}
     \begin{mat}
       \omega_0^2-\omega^2    &-\epsilon/2m  \\
       \epsilon/2I &\omega_0^2-\omega^2 
     \end{mat}
     \colvec{
       A_1 \\
       A_2}
     =
     \colvec{
       0 \\
       0}
   \end{equation*}
   我们要求使该方程组有非平凡
   解的频率~$\omega$。
   \begin{equation*}
      \begin{vmatrix}
       \omega_0^2-\omega^2    &-\epsilon/2m  \\
       \epsilon/2I           &\omega^2-\omega_0^2 
      \end{vmatrix}
      =0
    \end{equation*}
    这与专题正文中算过的行列式相同，
    只是第二列乘了 $-1$。
    用一个标量乘某一行或某一列，就相当于用该标量乘
    整个行列式。
    所以这个行列式是专题正文中那个行列式的相反数。
    但我们要令它为零，所以这无关紧要。
    根与专题正文中相同。
  \end{answer}
\item 用 Slinky Jr 弹簧和一个汤罐做一个 Wilberforce 摆。
  你可以在罐上钻孔装入螺栓，孔数两个或四个均可，
  用来调节罐的转动惯量，
  使竖直运动与转动运动的
  周期一致。
  \begin{answer}
    参见 \cite{Mewes}。
  \end{answer}
\end{exercises}
\index{Coupled Oscillators|)}
\endinput
% \end{document}
